\section{The connection on a uniformizing chart}
\label{sec:jet-connection}

Fix a framed point on the characteristic-zero fiber. In one copy of
a formal group chart let \(Z\) be the coordinates, \(F_0\) the
formal group law, \(A_a(Z)\) the action of \(a\in\mathcal O_E\),
and \(D_u\) the invariant derivation in tangent direction \(u\).
Thus \(D_u f(Z)\) is the coefficient of \(t\) in
\(f(F_0(Z,tu))\). All calculations with exponentials below take
place over a characteristic-zero field. The lattice on which
values are eventually bounded is the one in
Section~\ref{sec:local-theta-jets}.

\subsection{A local formula and its independence of gauge}

The rigidified cubical line provides multiplicative factors
\(\Phi_1,\Phi_2\) for the two partial addition laws of the
biextension, a symmetry factor, and an adjoint factor \(J_a\).
A canonical gauge is a series \(h(Z,W)\), zero on each axis, whose
partial additive coboundaries are the logarithms of the first two
factors. For instance,
\[
 \log\Phi_2(X,Y,Z)=h(X,Y)+h(X,Z)-h(X,F_0(Y,Z)).
\]
Its difference under interchanging the arguments is the logarithm
of the symmetry factor. The adjoint condition reads
\(\log J_a(X,Y)=h(A_aX,Y)-h(X,A_{\bar a}Y)\). With the element
\(w_0\) defining the primitive theta line, put
\(H_h(Z)=h(Z,A_{w_0}Z)\). If the germ in the primitive frame of a
section of the \(k\)-th power is \(f\), the chart evaluates a
monomial in the connection by
\[
      \left.D_{u_1}\cdots D_{u_j}\bigl(f\,\exp(kH_h)\bigr)
      \right|_{Z=0}.
\]
For multiple columns apply the corresponding derivatives in the
specified copies, with the product factor
\(\exp(k\sum_b H_h(Z_b))\).

We give the gauge construction to indicate where denominators enter.
Pass to the normalized formal logarithm of the commutative smooth
formal group. It is obtained by termwise integration of the
closed invariant differentials dual to the tangent basis. They are
closed since the invariant tangent fields commute, and their
constant matrix is the identity. The resulting coordinates turn
the group law into addition, because the pullback of each invariant
differential by addition is the sum of its two pullbacks. For a
symmetric, normalized additive cocycle \(C(x,y)\), allow inactive
coordinates as parameters, and put
\(g_j(x)=\partial_{y_j} C(x,0)\). Differentiating the cocycle equation
gives \(\partial_{y_j}C(x,y)=g_j(x+y)-g_j(y)\).
Symmetry, followed by taking the mixed partial derivatives, gives
\(\partial_i g_j=\partial_j g_i\). Choose the formal primitive \(h_1\)
of \(-\sum_jg_j\,dx_j\) with zero constant. It has zero linear part
in \(x\), and
\(C(x,y)=h_1(x)+h_1(y)-h_1(x+y)\), by differentiation in \(y\)
and evaluation at \(y=0\). Apply this to the first multiplication
law. After this change of frame, the interchange law for a
biextension makes the logarithm of the second defect additive in
its passive variable. Applying the same primitive in its active
variable preserves that additivity, since integration is linear in
the coefficients. This kills both defects.

The remaining endomorphism defects are biadditive. In their
\(\mathbb Q\)-vector space let \(l(a)\) act on the first argument
and \(r(a)\) by \(\bar a\) on the second. The logarithms \(\delta(a)\)
of the adjoint factors satisfy
\[
  \delta(ab)=l(b)\delta(a)+r(a)\delta(b).
\]
Indeed this is the logarithm of the composition law for the adjoint
trivializations, and \(l,r\) commute. In \(E\otimes_\mathbb Q E\)
take the separability element \(\sum_i x_i\otimes y_i\), characterized
by \(\sum_i x_i y_i=1\) and
\((a\otimes1)e=(1\otimes a)e\); it is the projector onto the
multiplication factor of the finite étale algebra \(E\otimes E\).
Extend \(l,r,\delta\) by rational linearity. Apply
\(x\otimes y\mapsto l(x)\delta(y)\) to the last identity to obtain
\(\delta(a)=l(a)v-r(a)v\), where \(v=\sum_i l(x_i)\delta(y_i)\).
Adding this biadditive correction kills the adjoint defects without
changing the multiplication laws. Finally, the residual logarithm
of the symmetry factor is alternating, biadditive, and commutes
with the adjoint actions. Subtracting half of it supplies the
last correction since \(h(Y,X)-h(X,Y)\) is the required symmetry
logarithm. This constructs \(h\).
\footnote{See \href{https://github.com/ChristianKudeba/alpoge-rh-verification/blob/alpoge-source-7cbf7a0c/src/Development/AbelianVarietyTheory/Proofs/FormalGroupCharts/CanonGaugeOfLog.lean}{\srcid{canonGauge_of_log}} and
\href{https://github.com/ChristianKudeba/alpoge-rh-verification/blob/alpoge-source-7cbf7a0c/src/Development/AbelianVarietyTheory/Proofs/FormalGroupCharts/DerivationInnerOfSeparabilityElement.lean}{the endomorphism correction}.}

For two choices \(h,h'\), their difference in logarithmic coordinates
is a bilinear form \(L\). In fact its two additive coboundaries
vanish, so differentiation over the characteristic-zero field makes
it linear in each argument. The adjoint law gives
\(L(au,v)=L(u,\bar a v)\). If \(u,v\) have eigencharacters
\(\tau,\sigma\) with \(\sigma\ne\bar\tau\), choose \(a\) with
\(\tau(a)\ne\sigma(\bar a)\). Subtraction gives \(L(u,v)=0\).
The diagonal change of gauge is a homogeneous quadratic polynomial
\(q(Z)=L(Z,dA_{w_0}Z)\) in logarithmic coordinates.

For a polynomial \(P\) in the derivative symbols, let
\(\langle P,S\rangle=P(\partial)S(0)\) in these coordinates.
Repeated differentiation of a monomial gives
\(\langle P,Z_iS\rangle=\langle\partial_iP,S\rangle\).
The letters in the source are tangent directions \(u_j\). Equivalently,
substitute \(Z=\sum t_j u_j\) in each column and use ordinary
directional differentiation in the \(t_j\); the same identity
applies even if the substitution is not an isomorphism. With
\(q_{\mathrm{tot}}=\sum_b q(Z_b)\), it follows that
\[
  \langle P,e^{kq_{\mathrm{tot}}}S\rangle
   =\sum_{j=0}^{\lfloor\deg P/2\rfloor}
       \frac{k^j}{j!}\,
          \langle q_{\mathrm{tot}}(\partial)^j P,S\rangle .
\]
The quadratic coefficients have support only in conjugate pairs.
The same form occurs on every column. Thus
\(q_{\mathrm{tot}}(\partial)\) is a linear combination of the
operators
\(\sum_b\partial_{(\tau,i,b)}\partial_{(\bar\tau,j,b)}\).
The defining harmonic equations annihilate \(P\) under each of them.
Only \(j=0\) contributes. Individual mixed derivatives may change
with the gauge; the asserted invariant is the harmonic combination.

For a change of formal coordinates, compose the two charts with
their normalized formal exponential in the marked tangent
directions. The resulting group homomorphisms agree since they
have the same differential and additive coordinates determine a
homomorphism in characteristic zero. A change of primitive frame
may initially give different frames at adjacent truncation orders.
Their ratio on the lower truncation is a unit normalized at the
axes. The equal cubical multiplication, symmetry and adjoint
factors imply that its logarithm in two arguments is biadditive,
symmetric and balanced, through that order. In logarithmic
coordinates it is a bilinear form through the truncation. Successive
subtractions of these bilinear forms make the frames compatible
to the desired order. Their cumulative diagonal change has the form
\(\exp(kq_{\mathrm{tot}})\) just treated. The total degree
\(D=\deg P\) suffices. Writing \(I=(Z)\), an invariant derivative
sends \(I^j\) into \(I^{j-1}\), by the Leibniz rule. Consequently
terms in \(I^{D+1}\) make no contribution to the evaluated
derivative. These facts give chart-independence at the exact order
used for the jet.
\footnote{\href{https://github.com/ChristianKudeba/alpoge-rh-verification/blob/alpoge-source-7cbf7a0c/src/Development/AbelianVarietyTheory/Proofs/FormalGroupCharts/CanonGaugeIndependenceHarmonic.lean}{Gauge change and conjugate eigendirections};
\href{https://github.com/ChristianKudeba/alpoge-rh-verification/blob/alpoge-source-7cbf7a0c/src/Formalism/SharpRestriction/Divisibility/GeneratorPowers/FrRatio.lean}{ratios of truncated frames};
\href{https://github.com/ChristianKudeba/alpoge-rh-verification/blob/alpoge-source-7cbf7a0c/src/Formalism/SharpRestriction/Divisibility/GeneratorPowers/CIndVerb.lean}{reduction of chart-independence to the finite jet}.}

\subsection{The normalized pullback factor}

Let \(\beta\bar\beta=N>0\) and write \(X=Z\) for the universal
formal point. The integral ratio between the normalized pullback
frame and the old frame is the following unit:
\[
 W(X)=
 \frac{J_\beta(X,A_{w_0\beta}X)}
 {\prod_{i=0}^{N-1}\Phi_2(X,A_{iw_0}X,A_{w_0}X)} .
\]
The normalization of the cubical line on the zero face supplies
constant one in this formula. To see the frame identity, identify
the primitive line at \(X\), divided by its origin fiber, with the
biextension line at \((X,A_{w_0}X)\).
The denominator composes the second partial addition
\(N\) times in the \(w_0\)-direction. The numerator transports the
\(\beta\)-action from the first to the second factor by Rosati,
where \(\bar\beta w_0\beta=Nw_0\). Both identifications agree at the
origin. Therefore a section with germ \(f\) in power \(k\) has
normalized pullback germ \(f'=W^k(f\circ A_\beta)\) in power \(kN\).

The logarithmic formula makes the entire adjustment explicit.
The denominator telescopes to
\(NH_h(X)-h(X,A_{Nw_0}X)\), since \(h(X,0)=0\).
The logarithm of its numerator is
\(H_h(A_\beta X)-h(X,A_{Nw_0}X)\). Hence
\[
      W=\exp\{H_h\circ A_\beta-NH_h\},\qquad
      f'e^{kNH_h}=(f e^{kH_h})\circ A_\beta .
\]
Applying an invariant derivative in direction \(u_{\tau,i}\) to a
pullback by \(A_\beta\) multiplies it by \(\tau(\beta)\). This follows
by differentiating the group-homomorphism law on the dual numbers.
Iteration, multiplication over the columns, and the fixed row
degrees of \(P\) give the character
\[
      \prod_{\tau,i}\tau(\beta)^{a_{\tau,i}} .
\]
All the harmonic monomials have that same factor.
\footnote{\href{https://github.com/ChristianKudeba/alpoge-rh-verification/blob/alpoge-source-7cbf7a0c/src/Development/AbelianVarietyTheory/Proofs/JetHarmonics/ThreePointSeriesExponentialOfDefect.lean}{Normalized cubical and adjoint factors};
\href{https://github.com/ChristianKudeba/alpoge-rh-verification/blob/alpoge-source-7cbf7a0c/src/Development/AbelianVarietyTheory/Proofs/JetHarmonics/HarmonicValueNormalisedPullbackChainRule.lean}{pullback germs and the character on the row degrees}.}

\subsection{Comparison on the uniformized frame}

Take the integral chart over \(T\), the framed generic point \(y\),
its polarized uniformization, and a generating frame on the jets
compatible with the rigidified primitive frame. Let \(F\) be a
characteristic-zero extension field of \(\operatorname{Frac}T\).
Suppose the uniformized sections \(\sigma_b\) correspond to the
geometric sections \(s_b\) under the specified generic-fiber
identification. Write \(f_b\) for their uniformized germs. To compare
the two evaluations it suffices, on each truncated jet, to make the
same base change in the coordinates, the tangent directions, the
cubical line and its generating frame. Indeed the two sections
restrict to the same section of the identified line. Dividing by
the common power of its generating frame gives the same truncated
series; equality may be checked on a trivializing open cover.
The canonical gauge also pulls back, since its identities are
identities of series and the homomorphism is between
characteristic-zero fields. The chart-independence just proved
identifies the scalar on these germs with the scalar defining the
connection on the original framed point. Finally, derivatives in
different columns act on independent variables, so evaluating the
box product of the \(f_b\) gives precisely the sum of products of
one-column values used to define the theta rule. This proves the
equality for the tied sections at any power \(k\) and harmonic \(P\).
\footnote{\href{https://github.com/ChristianKudeba/alpoge-rh-verification/blob/alpoge-source-7cbf7a0c/src/Development/AbelianVarietyTheory/Proofs/JetHarmonics/HarmonicValueConnectionIntegralChart.lean}{\srcid{harmonic_value_of_tied_sections}}.}

This field-valued comparison can also be used at a finite residue
prime. An integral assertion about the connection requires additional
control. Here is the unit hypothesis in one such assertion. Let
\(u,u'\) be integral directions of character \(\sigma\). Let \(B_h\)
be the matrix of the bidegree-\((1,1)\) part of \(h\). Taking that
component of the adjoint equation shows that
\[
       (dA_a)^t B_h-B_h dA_{\bar a}
\]
has integral entries. The logarithm of the integral adjoint unit
has the same bidegree-\((1,1)\) coefficients as the unit itself:
its rigidification makes it one on both axes, so higher powers of
its augmentation have bidegree at least \((2,2)\). Contracting the
matrix between \(u^t\) and \(u'\) proves that
\((\sigma(a)-\sigma(\bar a))u^tB_hu'\) is integral.
Thus an \(a\) with \(\sigma(a)-\sigma(\bar a)\in T^\times\) gives
integrality of this contraction. The additional operator \(dA_{w_0}\)
on \(u'\) multiplies it by the integral eigenvalue \(\sigma(w_0)\).
This unit-separation hypothesis belongs to this coefficient argument.
The separate projector for integral theta translates in
Section~\ref{sec:theta-translate-span} uses \(q\in T^\times\).
On the horizontal boundary both arguments are in residue
characteristic zero, where a nonzero element of the coefficient
field has valuation zero under the structural map through \(T\).
\footnote{\href{https://github.com/ChristianKudeba/alpoge-rh-verification/blob/alpoge-source-7cbf7a0c/src/Development/AbelianVarietyTheory/Proofs/JetHarmonics/CanonicalGaugeEigenvectorPairingIntegral.lean}{The precise same-character integrality lemma}.}

\subsection{Why the boundary weight shift has zero cost}

Let \(e_Z\) be the idempotent of the component. If a cusp section
\(f=e_Z f\) has nonzero evaluation, then \(e_Z(y)=1\).
The source constructs a distinguished unit section \(K_4\) and a
unit rule \(\rho\) with \(e_Z\rho^4=a\,e_ZK_4\),
\(a\) a nonzero coefficient-field element depending on \(Z\).
The integral-model theorem for the theta scalar gives
\(v_T(K_4(y))=0\) here. The map from the stalk of the horizontal
boundary sends the coefficient field into \(T\). Applying this to
both \(a\) and \(a^{-1}\) gives \(v_T(a)=0\). Evaluating the fourth
power identity gives \(4v_T(\rho(y))=0\), hence the weight-one
twisting section \(t=e_Z\rho^{-1}\) has order zero. Its inverse on
the component is \(e_Z\rho\). The input on parallel weights gives
another invertible twisting section \(t_{-w}\) of weight
\(\operatorname{par}(-w)\), of order zero at the same point.
Its boundedness and that of its inverse preserve the cusp condition;
the same property is part of the construction of \(t\).
Thus \(f\,t_{-w}t^{-m}\) removes both \(\operatorname{par}(w)\) and
the scalar weight \(m\), preserving nonvanishing and the exact order.
The boundary theorem is stated for \(4\mid m\); in this application
\(m=n/2\) and \(8\mid n\). The final multiplication uses only the
integer power once the weight-one twist is in hand. The fourth
power relation enters in constructing and controlling that twist.

Section~\ref{sec:key-different} derives the global fourth root from
the complex theta normalization and gives the finite-place different
bound. The model unit theorem for \(K_4\) and the construction of the
parallel-weight twist also enter the boundary argument.
\footnote{\href{https://github.com/ChristianKudeba/alpoge-rh-verification/blob/alpoge-source-7cbf7a0c/src/Development/AbelianVarietyTheory/Proofs/DivisorProofs/IntegralFramePointStalk.lean}{Construction of the weight-one twist and its valuation};
\href{https://github.com/ChristianKudeba/alpoge-rh-verification/blob/alpoge-source-7cbf7a0c/src/Development/AutomorphicArithmetic/ThetaCorrespondence/ThetaPairLifts/ThirdPlaceSeries/BoundaryValuationTwistInvariant.lean}{the boundary twist at the prescribed point}.}
